\section{Задача без номера}

Пусть $\ee = (e_1, e_2, e_3, e_4) \subset V$ --- ортонормированный базис.

\subsection{Угол между сектором и плоскостью}

Найдите угол между вектором~$v=\ee(1, -5, -3, -5)^\top$ и плоскостью~$\pi$, заданной системой уравнений
\begin{equation}
	\label{eq:num:b:plane_system}
	\begin{system}
		-4x + 3y + 2z + t 
		& = 0,
		\\
		x - 2y + z
		& = 0.
	\end{system}
\end{equation}

\begin{solution}
	В нашем случае плоскость задана как ядро матрицы
	\[
		B :=
		\begin{pmatrix}
			-4 & 3 & 2 & 1 \\
			1 & -1 & 1 & 0
		\end{pmatrix}
		.
	\]
	Чтобы это увидеть, заметим, что система~\eqref{eq:num:b:plane_system} эквивалентна уравнению~$Bw = 0$, где $w = (x, y, z, t)^\top$.

	Если бы плоскость~$\pi$ была задана как двумерная плоскость в трёхмерном пространстве, то найти косинус угла~$\theta$ между~$v$ и~$\pi$ можно было бы, найдя скалярное произедение~$\spr{v}{n}$, где~$n$ --- нормаль к~$\pi$. Тогда
	\[
		\cos{\theta} =
		\frac{\spt{v}{n}}{\norm{v} \cdot \norm{n}}
		.
	\]
	Мы воспользовались бы тем, что ортогональное дополнение до~$\pi$ одномерно.

	Но в нашем случае ортогональное дополнение до плоскости~$\pi$ двумерно, как и сама плоскость~$\pi$, и поэтому надо искать угол между одномерным подпространством~$\anspan{\vec{v}}$ и двумерным подпространством~$\pi$ либо его ортогональным дополнением~$\pi^{\bottom}$.
\end{solution}

